\documentclass{article}

\usepackage{amsmath,amssymb}
\usepackage{xurl}
\usepackage[hidelinks]{hyperref}
\setlength{\emergencystretch}{2em}

\input{command}

\title{Slater Strong Duality: Finiteness Conditions for Attainment}
\author{}
\date{2026-08-24}

\begin{document}
\maketitle
\gapnote{false-source}{resolved}

\section{Source passage}

Let $V$ and $W$ be finite-dimensional real inner-product spaces, let
$K\subseteq V$ be a closed pointed convex cone with nonempty interior, and let
$T:V\to W$ be linear. For $c\in V$ and $b\in W$, Chapter~4 of
Ref.~\cite{Wolf2012Quantum} defines the primal and dual values by
\begin{align}
  C_p
    &=\inf\{\langle c|x\rangle:x\in K,\ T(x)=b\},
  \label{eq:wolf_slater_attainment_conditions_primal_value}\\
  C_d
    &=\sup\{\langle b|y\rangle:c-T^*(y)\in K^*\}.
  \label{eq:wolf_slater_attainment_conditions_dual_value}
\end{align}
At lines~72--78 of the local source, primal strict feasibility is said to imply
$C_p=C_d$ and attainment of the dual supremum, while dual strict feasibility
is said to imply attainment of the primal infimum
\cite[Chapter~4]{Wolf2012Quantum}.

The semidefinite specialization at lines~100--105 of the same source contains
the same missing finiteness condition and a separate directional error
\cite[Chapter~4]{Wolf2012Quantum}. It assumes either a positive-definite primal
matrix $X$ satisfying $\tr(F_iX)=b_i$, or a dual vector $y$ whose slack
$F_0-\sum_i y_iF_i$ is positive definite, and then states in both cases that
the dual extremum is attained. Primal strict feasibility yields dual
attainment when the primal value is finite. Dual strict feasibility instead
yields primal attainment when the dual value is finite; it does not generally
yield dual attainment.

\section{Missing boundary conditions}

Strict feasibility on one side does not ensure that the opposite feasible set
is nonempty or that the relevant optimum is finite. Both failures occur for
$V=\mathbb R$, $K=\mathbb R_{\geq0}$, and $T=0$.

First, set $b=0$ and $c=-1$. The point $x=1$ is strictly primal feasible, but
\begin{align}
  \inf_{x\geq0}(-x)
    &=-\infty.
  \label{eq:wolf_slater_attainment_conditions_unbounded_primal}
\end{align}
Dual feasibility would require
\begin{align}
  -1
    &\in\mathbb R_{\geq0},
  \label{eq:wolf_slater_attainment_conditions_empty_dual_requirement}
\end{align}
which is impossible. Thus the dual feasible set is empty and has no optimizer.

Second, set $b=1$ and $c=1$. The dual slack is $1$ for every $y$, so the dual
problem is strictly feasible. The primal equality constraint becomes
\begin{align}
  0
    &=1,
  \label{eq:wolf_slater_attainment_conditions_infeasible_primal_constraint}
\end{align}
so the primal feasible set is empty and has no optimizer.

These examples are also semidefinite programs on $1\times1$ Hermitian
matrices with one equality constraint. For the first example, take
$F_0=-1$, $F_1=0$, and $b_1=0$. The matrix $X=1$ is positive definite and
satisfies $\tr(F_1X)=b_1$, but the primal objective is unbounded below and the
dual slack satisfies
\begin{align}
  F_0-y_1F_1
    &=-1
  \label{eq:wolf_slater_attainment_conditions_negative_dual_slack}
\end{align}
for every $y_1$, so it is never positive semidefinite. For the second example,
take $F_0=1$, $F_1=0$, and $b_1=1$. Every dual slack is the positive-definite
matrix $1$, while the primal equality
\begin{align}
  \tr(F_1X)
    &=1
  \label{eq:wolf_slater_attainment_conditions_impossible_sdp_equality}
\end{align}
is impossible.

Therefore, the attainment statements at lines~72--78 and their semidefinite
specialization at lines~100--105 require a finiteness condition, or an
equivalent condition combining opposite-side feasibility and boundedness
\cite[Chapter~4]{Wolf2012Quantum}.

A finite-valued semidefinite program isolates the directional error at
line~105. In the convention of Ref.~\cite{Wolf2012Quantum}, set
\begin{align}
  F_0
    &=\begin{pmatrix}0&1\\1&0\end{pmatrix},
  \label{eq:wolf_slater_attainment_conditions_directional_F0}\\
  F_1
    &=\begin{pmatrix}-1&0\\0&0\end{pmatrix},
  \label{eq:wolf_slater_attainment_conditions_directional_F1}\\
  F_2
    &=\begin{pmatrix}0&0\\0&-1\end{pmatrix},
  \label{eq:wolf_slater_attainment_conditions_directional_F2}\\
  b
    &=(-1,0).
  \label{eq:wolf_slater_attainment_conditions_directional_b}
\end{align}
Writing the dual variable as $y=(x,z)$ gives
\begin{align}
  F_0-xF_1-zF_2
    &=\begin{pmatrix}x&1\\1&z\end{pmatrix},
  \label{eq:wolf_slater_attainment_conditions_directional_slack}\\
  \langle b|y\rangle
    &=-x.
  \label{eq:wolf_slater_attainment_conditions_directional_objective}
\end{align}
The dual problem therefore maximizes $-x$ subject to positivity of the matrix
in (\ref{eq:wolf_slater_attainment_conditions_directional_slack}). It is
strictly feasible at $(x,z)=(2,2)$. Feasibility implies
\begin{align}
  x
    &>0,
  \label{eq:wolf_slater_attainment_conditions_directional_x_positive}\\
  z
    &>0,
  \label{eq:wolf_slater_attainment_conditions_directional_z_positive}\\
  xz
    &\geq1.
  \label{eq:wolf_slater_attainment_conditions_directional_determinant}
\end{align}
For every $t>0$, the point $(x,z)=(t,t^{-1})$ is feasible. Hence
\begin{align}
  \sup\{-x:(x,z)\text{ is dual feasible}\}
    &=0,
  \label{eq:wolf_slater_attainment_conditions_unattained_dual_supremum}
\end{align}
but the supremum is approached only as $t\downarrow0$ and is not attained.

The associated primal constraints are
\begin{align}
  -X_{11}
    &=-1,
  \label{eq:wolf_slater_attainment_conditions_primal_first_diagonal}\\
  -X_{22}
    &=0.
  \label{eq:wolf_slater_attainment_conditions_primal_second_diagonal}
\end{align}
Positive semidefiniteness then forces
\begin{align}
  X
    &=\begin{pmatrix}1&0\\0&0\end{pmatrix},
  \label{eq:wolf_slater_attainment_conditions_primal_optimizer}
\end{align}
and its objective value is
\begin{align}
  \tr(F_0X)
    &=0.
  \label{eq:wolf_slater_attainment_conditions_primal_optimum}
\end{align}
Thus the primal infimum is attained and equals the dual supremum, exactly as
the correctly oriented Slater theorem asserts.

\section{Corrected statement}

In the signs and adjoint convention of Ref.~\cite{Wolf2012Quantum}, the
finite-dimensional Slater theorem has the following two directions:
\begin{align}
  \text{primal strict feasibility and }C_p\in\mathbb R
    &\Longrightarrow
      C_p=C_d\text{ and dual attainment},
  \label{eq:wolf_slater_attainment_conditions_corrected_primal_direction}\\
  \text{dual strict feasibility and }C_d\in\mathbb R
    &\Longrightarrow
      C_p=C_d\text{ and primal attainment}.
  \label{eq:wolf_slater_attainment_conditions_corrected_dual_direction}
\end{align}

The formal development proves both implications in
\doclink{QICLean/Analysis/ConicProgram}.\footnote{The primal-strict direction
is given by
\leanid{ConicProgram.exists_isDualOptimizer_of_isPrimalStrictlyFeasible_of_primalValue_eq_coe}
and
\leanid{ConicProgram.values_eq_of_isPrimalStrictlyFeasible_of_primalValue_eq_coe}.
The dual-strict direction is given by
\leanid{ConicProgram.exists_isPrimalOptimizer_of_isDualStrictlyFeasible_of_dualValue_eq_coe}
and
\leanid{ConicProgram.values_eq_of_isDualStrictlyFeasible_of_dualValue_eq_coe}.}
The first pair assumes primal strict feasibility and $C_p=p\in\mathbb R$; the
second assumes dual strict feasibility and $C_d=p\in\mathbb R$. The
declarations
\leanid{ConicProgram.primalValue_eq_coe_iff_isGLB} and
\leanid{ConicProgram.dualValue_eq_coe_iff_isLUB} identify these hypotheses
with the corresponding real greatest-lower-bound and least-upper-bound
conditions whenever the relevant feasible set is nonempty.

The direct semidefinite specializations are proved in
\doclink{QICLean/Analysis/SemidefiniteDuality}.\footnote{The primal-strict
declarations are
\leanid{SemidefiniteProgram.exists_dualOptimizer_of_primalStrict_of_primalValue_eq_coe}
and
\leanid{SemidefiniteProgram.values_eq_of_primalStrict_of_primalValue_eq_coe}.
The dual-strict declarations are
\leanid{SemidefiniteProgram.exists_primalOptimizer_of_dualStrict_of_dualValue_eq_coe}
and
\leanid{SemidefiniteProgram.values_eq_of_dualStrict_of_dualValue_eq_coe}.}
A positive-definite primal point together with $C_p=p\in\mathbb R$ gives dual
attainment and value equality. A positive-definite dual slack together with
$C_d=p\in\mathbb R$ gives primal attainment and value equality.

\section{Separation argument}

Let $C$ be a cone, let $A$ be an equality map, let $d$ be an objective vector,
and let $r$ be the right-hand side. The proof separates the lifted epigraph
\begin{align}
  D
    &=\{(A(x),\langle d|x\rangle+s):x\in C,\ s\geq0\}
      \subseteq\operatorname{ran}(A)\times\mathbb R.
  \label{eq:wolf_slater_attainment_conditions_lifted_epigraph}
\end{align}
Restricting the codomain to $\operatorname{ran}(A)$ makes the lifted linear
map surjective and therefore open in finite dimension. Strict feasibility
gives $D$ nonempty interior. If $p$ is the finite greatest lower bound, then
\begin{align}
  (r,p)
    &\in\overline D\setminus\operatorname{int}D.
  \label{eq:wolf_slater_attainment_conditions_boundary_point}
\end{align}
A supporting functional $f$ therefore satisfies
\begin{align}
  f(r,p)
    &=0,
  \label{eq:wolf_slater_attainment_conditions_supporting_vanishes}\\
  f(z)
    &\leq0
    \qquad\text{for every }z\in D.
  \label{eq:wolf_slater_attainment_conditions_supporting_nonpositive}
\end{align}
Its vertical coefficient
\begin{align}
  \alpha
    &=f(0,1)
  \label{eq:wolf_slater_attainment_conditions_vertical_coefficient}
\end{align}
is strictly negative. If $\alpha=0$, strict feasibility and openness of the
nonzero functional would make $0$ an interior point of $(-\infty,0]$, which
is impossible. Normalization by $-\alpha$ gives the required dual multiplier
and its attained objective value.

For the converse direction, write the dual problem as a primal problem on
$W\times K^*$ with equality map and objective
\begin{align}
  (y,z)
    &\longmapsto T^*(y)+z,
  \label{eq:wolf_slater_attainment_conditions_converse_equality_map}\\
  (y,z)
    &\longmapsto\langle-b|y\rangle.
  \label{eq:wolf_slater_attainment_conditions_converse_objective}
\end{align}
The multiplier returned by the same separation argument satisfies
\begin{align}
  T(-x)
    &=b,
  \label{eq:wolf_slater_attainment_conditions_converse_constraint}\\
  -x
    &\in K^{**}=K.
  \label{eq:wolf_slater_attainment_conditions_converse_cone_membership}
\end{align}
These are precisely the primal feasibility conditions in Wolf's signs. The
argument never assumes that a linear image of a closed cone is closed. It uses
only the closure membership in
(\ref{eq:wolf_slater_attainment_conditions_boundary_point}), which follows
from the greatest-lower-bound property.

\section{Verdict}

\textbf{Verdict: false as printed; corrected theorem resolved.} The two
one-dimensional examples show that neither the general assertions at
lines~72--78 nor their semidefinite specialization at lines~100--105 follows
from strict feasibility alone. The $2\times2$ example also shows that dual
strict feasibility and a finite dual supremum do not imply dual attainment, so
line~105 names the wrong optimizer in the dual-strict case. With the missing
finite-real optimum hypothesis, value equality and attainment on the opposite
side are now formalized by separate declarations. This resolves
\href{https://github.com/LionSR/QICLean/issues/110}{QICLean issue \#110} while
retaining the source correction. The same correction governs the semidefinite
specialization tracked by
\href{https://github.com/LionSR/QICLean/issues/111}{QICLean issue \#111}.

\bibliographystyle{alpha}
\bibliography{references}

\end{document}
